\begin{blocksection}
\question A rock-paper-scissors app stores its state in a model with three fields: \lstinline$your_score$ (an \lstinline$Int$), \lstinline$computer_score$ (an \lstinline$Int$), and \lstinline$rounds$ (a list of every round played so far, each recording both players' choices). A new feature removes the most recent rounds from \lstinline$rounds$. What must also happen?

\begin{choices}
\choice Nothing, because scores are not part of the model's state
\choice The scores must be recomputed from the remaining rounds, because they can be derived from \lstinline$rounds$ and could otherwise disagree with it
\choice \lstinline$view$ must be changed to ignore the scores
\choice The scores must be removed from the model, because stored values that can be derived from other fields are not allowed
\end{choices}

\begin{solution}
(B). The scores can be computed from \lstinline$rounds$, so after rounds are removed, the stored scores no longer match the history unless they are recomputed. Storing derived values is convenient, but every change to \lstinline$rounds$ must update them too.
(A) is wrong because the scores are stored in the model.
(C) doesn't help, since \lstinline$view$ just displays whatever the model holds.
(D) overcorrects, storing a derived value is allowed, it just means you must keep it consistent.
\end{solution}
\end{blocksection}

\begin{blocksection}
\question The same app lets the user step back through past rounds. The model needs to remember which round is currently displayed. Which design is least likely to cause bugs?

\begin{choices}
\choice Add a field \lstinline$displayed: Round$ holding a copy of the round being shown
\choice Add a field \lstinline$back: Int$, the number of rounds before the latest one, and find the displayed round from \lstinline$rounds$ and \lstinline$back$
\choice Add both \lstinline$displayed: Round$ and \lstinline$back: Int$, so \lstinline$view$ can use whichever is convenient
\choice Remove \lstinline$rounds$ and store only \lstinline$displayed$
\end{choices}

\begin{solution}
(B). \lstinline$rounds$ already holds every round. \lstinline$back$ is just a position in that list, so the displayed round can always be found from \lstinline$rounds$ and \lstinline$back$.
(A) and (C) store a copy of a round, which can become wrong when \lstinline$rounds$ changes.
(D) deletes the history, so there is nothing to step back through.
\end{solution}
\end{blocksection}

\begin{blocksection}
\question In the same app, \lstinline$view$ shows both scores and the hands from the most recent round in \lstinline$rounds$. A developer writes the \lstinline$UserChose$ case of \lstinline$update$ so that it correctly increases the scores but forgets to add the new round to \lstinline$rounds$. The user clicks Rock. What does the page show?

\begin{choices}
\choice The updated scores, but the hands from the previous round (or no hands, if this was the first click)
\choice The old scores and the old hands, because nothing changed
\choice The updated scores and the hands from the round the user just played, because Lustre remembers the click
\choice An error, because the model's fields are inconsistent
\end{choices}

\begin{solution}
(A). \lstinline$view$ only sees the model that \lstinline$update$ returns. The scores were updated, so the page shows the new scores. The new round was never added to \lstinline$rounds$, so \lstinline$view$ still finds the previous round at the front of the list (or an empty list on the first click, which shows no hands).
(B) is wrong because the scores did change.
(C) is wrong because the click isn't stored anywhere except in the model, so Lustre can't show a round the model doesn't contain.
(D) is wrong because the model's types are still valid, so nothing crashes; the page just shows stale hands.
\end{solution}
\end{blocksection}

\begin{blocksection}
\question What does \lstinline$element.to_string$ return for this element?

\begin{lstlisting}
html.button([attribute.class("button"), event.on_click(UserClickedReset)], [
  html.text("Start over"),
])
\end{lstlisting}

\begin{choices}
\choice \lstinline$<button class="button" onclick="UserClickedReset">Start over</button>$
\choice \lstinline$<button class="button">Start over</button>$
\choice \lstinline$<button class="button">UserClickedReset</button>$
\choice \lstinline$<button>Start over</button>$
\end{choices}

\begin{solution}
(B). \lstinline$element.to_string$ produces HTML text, and event handlers like \lstinline$event.on_click$ are associated with nodes in the tree but are not part of the HTML text. So only the class attribute and the text child appear.
(A) wrongly includes the event as an \lstinline$onclick$ attribute.
(C) confuses the message with the button's text child.
(D) leaves out the \lstinline$class$ attribute, which does appear in the HTML.
\end{solution}
\end{blocksection}